\documentclass[11pt]{article}
\newcommand{\HomeworkHeader}{习题二}
% Shared LaTeX preamble for homework/*/main.tex

% --- Base layout ---
\usepackage[a4paper, top=2.5cm, bottom=2.5cm, left=2cm, right=2cm]{geometry}
\usepackage{fontspec}
\usepackage[UTF8]{ctex}%latex支持中文宏包
\usepackage{amsmath, amssymb, amsthm}
\usepackage{mathtools}
\usepackage[svgnames]{xcolor}
\usepackage{pgfplots}
\pgfplotsset{compat=1.18}
\usepackage[most]{tcolorbox}
\usepackage{enumitem}
% Modified: Change label from hyphen (-) to bullet point ($\bullet$) for better visibility
\setlist[itemize]{label=$\bullet$}
\setlist[enumerate]{label=(\arabic*)}


% --- Colors ---
\definecolor{primaryColor}{RGB}{46, 52, 64}
\definecolor{accentColor}{RGB}{94, 129, 172}
\definecolor{boxFill}{RGB}{236, 240, 241}
\definecolor{solLine}{RGB}{163, 190, 140}
\definecolor{expandColor}{RGB}{191, 97, 106}

% --- TikZ styles (DFA / NFA / PDA) ---
\usepackage{tikz}
\usetikzlibrary{automata, positioning, arrows.meta, bending, backgrounds}
\tikzset{
    dfa_state/.style={
        state,
        thick,
        draw=accentColor,
        fill=boxFill,
        text=primaryColor,
        minimum size=1.1cm
    },
    dfa_edge/.style={
        ->,
        >=stealth,
        thick,
        draw=primaryColor,
        auto
    },
    pda_node/.style={
        state,
        thick,
        draw=accentColor,
        fill=boxFill,
        text=primaryColor,
        minimum size=1.2cm
    },
    pda_edge/.style={
        ->,
        >=stealth,
        thick,
        draw=primaryColor,
        auto
    },
    rule_expand/.style={
        pda_edge,
        draw=expandColor,
        text=expandColor
    },
    rule_match/.style={
        pda_edge,
        draw=solLine,
        text=solLine!80!black
    }
}

% --- Header / footer ---
\usepackage{fancyhdr}
\pagestyle{fancy}
\setlength{\headheight}{13.6pt}%避免页眉高度警告
\fancyhf{}
\fancyhead[L]{\kaishu 中国科学院大学}
\fancyhead[C]{林俊杰2023K8009916012}
\fancyhead[R]{\kaishu 理论计算机}
\usepackage{lastpage}
\cfoot{\small 第 \thepage \ 页，共 \pageref{LastPage} 页}

% --- Problem box ---
\newtcolorbox[use counter=problemSeq]{problem}[1][]{
    enhanced,
    breakable,
    colback=boxFill,
    colframe=accentColor,
    coltitle=white,
    fonttitle=\bfseries\large,
    title={题目 ~\theproblemSeq \quad #1},
    boxrule=0.5mm,
    arc=3mm,
    drop shadow=black!15!white,
    attach boxed title to top left={xshift=0.5cm, yshift*=-3mm},
    boxed title style={colback=accentColor},
    % 关键修改：每当创建一个新问题盒子时，自动将步骤计数器重置为0
    code={\setcounter{stepcount}{0}}
}

% --- Solution 环境 (红色盒子样式，支持可选序号) ---
% 修改说明：
% 1. [1][] 表示接受一个可选参数，默认为空
% 2. title={...} 中使用 \ifstrempty 判断是否有参数，如果有则添加空格和参数
\newtcolorbox{solution}[1][]{
    enhanced, breakable,
    colback=red!3!white,        % 背景色：淡粉
    colframe=red!75!black,      % 边框色：深红
    coltitle=white,
    fonttitle=\bfseries\Large\itshape,
    title={Solution\ifstrempty{#1}{}{~#1}}, % 动态标题：Solution (1)
    boxrule=0.5mm,
    arc=3mm,
    drop shadow=black!15!white,
    attach boxed title to top left={xshift=0.5cm, yshift*=-3mm},
    boxed title style={colback=red!75!black},
    before skip=0.5em,          % 盒子前的间距
    after skip=2em,             % 盒子后的间距
    before upper={\setcounter{stepcount}{0}} % 每个解答开始时重置步骤计数
}

% --- Title ---
\newcommand{\makeCustomTitle}[1]{
    \begin{center}
        \vspace*{1cm}
        {\Huge \bfseries \color{primaryColor} 作业 #1}\\
        \vspace{0.5cm}
        {\Large \color{accentColor} \today}
        \vspace{1.2cm}
    \end{center}
}

% --- 自定义环境定义 ---

% --- 自定义计数器与环境 ---
\newcounter{stepcount}
\newcounter{casecount}[stepcount]
\newcounter{problemSeq} % 新增：专门用于 Problem 的独立计数器

% 【关键修改】挂载钩子：每当遇到 \section 命令（含 \section*），将 problemSeq 计数器重置为 0
\pretocmd{\section}{\setcounter{problemSeq}{0}}{}{}

% 2. 定义带自动编号的步骤环境
\newenvironment{proofstep}[1]
  {%
    \refstepcounter{stepcount}% 增加计数
    \par\vspace{1.5em}%
    \noindent\textbf{\large \thestepcount. #1}% 输出 "1. 标题"
    \par\vspace{0.5em}%
  }
  {\par}

% 3. 定义无编号的结论环境 (用于最后的结论，格式同步骤)
\newenvironment{proofresult}[1]
  {%
    \par\vspace{1.5em}%
    \noindent\textbf{\large #1}% 仅输出标题
    \par\vspace{0.5em}%
  }
  {\par}

% 4. 定义带自动编号的情形环境
\newenvironment{proofcase}[1]
  {%
    \refstepcounter{casecount}% 增加计数
    \par\vspace{1em}%
    \noindent\textbf{情形 \thecasecount：#1}% 输出 "情形 1：标题"
    \par\vspace{0.2em}%
  }
  {\par}
\begin{document}
\makeCustomTitle{习题二}
% 红框选题：7（2）、9、11、14、15、20（全部三问）、22。
% 来源 assets/习题二.pdf：第1页 7（2）；第2页 9、11、14、15；第3页 20；第4页 22。
% 对应教材第64--67页。

\begin{problem}[原题 7（2）]
确定下列系统完全能控时参数 $a,b$ 的取值范围：
\[
\dot x=Ax+Bu,\qquad
A=\begin{bmatrix}-2&1&0\\0&-2&0\\0&0&-2\end{bmatrix},\quad
B=\begin{bmatrix}a&1\\2&4\\b&1\end{bmatrix}.
\]
\end{problem}
\begin{solution}
\begin{proofstep}{把能控性问题化为列向量的秩问题}
三阶系统完全能控的充要条件是
\[
\operatorname{rank}\mathcal C(A,B)=3,\qquad
\mathcal C(A,B)=[B\ AB\ A^2B].
\]
本题 $A$ 与标量矩阵 $-2I$ 只差一个非零元素，因此先写成
\[
A=-2I+N,\qquad N=\begin{bmatrix}0&1&0\\0&0&0\\0&0&0\end{bmatrix},\quad N^2=0.
\]
由此得 $AB=-2B+NB$、$A^2B=4B-4NB$。另一方面，
$NB=AB+2B$，所以 $[B\ AB\ A^2B]$ 与 $[B\ NB]$ 的列空间相同。
这样只需判断 $[B\ NB]$ 的秩。
\end{proofstep}

\begin{proofstep}{利用第一坐标方向化简矩阵}
左乘 $N$ 的作用是把 $B$ 的第二行放到第一行，其余行置零，故
\[
NB=\begin{bmatrix}2&4\\0&0\\0&0\end{bmatrix}.
\]
其中第一列为 $2e_1$、第二列为 $4e_1$，因此 $e_1$ 始终属于能控子空间。
从 $B$ 的第一列减去 $ae_1$、第二列减去 $e_1$，均不改变与 $e_1$ 一起张成的空间，得到
\[
\operatorname{rank}[B\ NB]
=\operatorname{rank}\begin{bmatrix}1&0&0\\0&2&4\\0&b&1\end{bmatrix}
=1+\operatorname{rank}\begin{bmatrix}2&4\\b&1\end{bmatrix}.
\]
这也说明参数 $a$ 只影响已经存在的第一坐标方向，不影响能控秩。
\end{proofstep}

\begin{proofstep}{区分参数的一般取值和例外取值}
剩余二阶矩阵的行列式为
\[
\det\begin{bmatrix}2&4\\b&1\end{bmatrix}=2-4b.
\]
若 $b\ne\tfrac12$，该矩阵秩为2，因此原能控秩为3；若 $b=\tfrac12$，
第二行恰为第一行的 $\tfrac14$，而第一行非零，因此其秩为1，原能控秩恰为2。
于是完整的取值范围为
\[
\boxed{a\in\mathbb R\ \text{任意},\qquad b\in\mathbb R\setminus\{\tfrac12\}.}
\]
\end{proofstep}
\end{solution}

\begin{problem}[原题 9]
设单输入系统 $(A,b)$ 满足
\[
A=\begin{bmatrix}A_1&0\\0&A_2\end{bmatrix},\qquad
b=\begin{bmatrix}b_1\\b_2\end{bmatrix},\qquad
A_i\in\mathbb R^{n_i\times n_i},\quad b_i\in\mathbb R^{n_i}.
\]
证明：$(A,b)$ 能控，当且仅当 $(A_1,b_1)$、$(A_2,b_2)$ 均能控，且 $A_1,A_2$ 无公共特征根。
\end{problem}
\begin{solution}
\begin{proofstep}{选择适合块对角结构的判据}
PBH 左特征向量判据指出：$(A,b)$ 能控，当且仅当对每个 $\lambda\in\mathbb C$，
满足 $qA=\lambda q$ 的非零行向量 $q$ 都有 $qb\ne0$。
虽然系统矩阵是实矩阵，特征值可能为复数，故在复数域使用此判据。

将 $q$ 按分块写成 $[q_1\ q_2]$，则
\[
qA=\lambda q
\iff q_1A_1=\lambda q_1,\quad q_2A_2=\lambda q_2,
\qquad qb=q_1b_1+q_2b_2.
\]
后面的必要性和充分性都由这两个关系推出。
\end{proofstep}

\begin{proofstep}{必要性：整体能控要求两个子系统都能控}
假定 $(A,b)$ 能控。若 $(A_1,b_1)$ 不能控，按 PBH 判据存在
$q_1\ne0$ 和 $\lambda$，使 $q_1A_1=\lambda q_1$、$q_1b_1=0$。
把该行向量补成 $q=[q_1\ 0]$，则
\[
q\ne0,\qquad qA=[q_1A_1\ 0]=\lambda q,\qquad qb=q_1b_1=0.
\]
这与整体能控矛盾，所以第一个子系统必须能控。对第二个子系统使用
$q=[0\ q_2]$ 完全同理。
\end{proofstep}

\pagebreak
\begin{proofstep}{必要性：两个子系统不能有公共特征根}
继续假定整体能控。若 $A_1,A_2$ 有公共特征根 $\lambda$，各取相应的非零左特征向量
$q_1,q_2$。上一阶段已证两个子系统均能控，所以
\[
\alpha=q_1b_1\ne0,\qquad \beta=q_2b_2\ne0.
\]
这里 $\alpha,\beta$ 是标量，关键就在于两个子系统共用同一个标量输入。
因此可以把它们的输入作用相互抵消，构造
\[
q=[\beta q_1\ -\alpha q_2]\ne0,\qquad
qA=\lambda q,\qquad qb=\beta\alpha-\alpha\beta=0,
\]
再次违反 PBH 判据。故两者无公共特征根。
\end{proofstep}

\begin{proofstep}{充分性：逐个检查整体系统的左特征向量}
现在假定两个子系统均能控，且两者无公共特征根。
任取 $q=[q_1\ q_2]\ne0$ 满足 $qA=\lambda q$。
若 $q_1,q_2$ 同时非零，则 $\lambda$ 同时是 $A_1,A_2$ 的特征值，矛盾。
因此恰有一个分块非零。

不妨 $q_1\ne0$。这时 $\lambda$ 是 $A_1$ 的特征值而不是 $A_2$ 的特征值，
$A_2-\lambda I$ 可逆，由 $q_2(A_2-\lambda I)=0$ 得 $q_2=0$。
第一个子系统能控又给出
\[
qb=q_1b_1\ne0.
\]
另一种情况同理。故不存在使 $qb=0$ 的非零左特征向量，整体系统能控。
必要性和充分性均已证明。
\end{proofstep}
\end{solution}

\begin{problem}[原题 11]
证明：系统 $(A,B)$ 能控的充要条件是，对每个 $n\times n$ 矩阵 $X$，若 $AX=XA$ 且 $XB=0$，则必有 $X=0$。
\end{problem}
\begin{solution}
\begin{proofstep}{必要性：由交换关系把零作用扩展到全部能控方向}
假定 $(A,B)$ 能控，且 $AX=XA$、$XB=0$。交换关系意味着
$XA^2=(XA)A=(AX)A=A(XA)=A^2X$；同样归纳得到
$XA^k=A^kX$。所以
\[
XA^kB=A^kXB=0\quad(k\ge0),\qquad
X\underbrace{[B\ AB\ \cdots\ A^{n-1}B]}_{\mathcal C(A,B)}=0.
\]
能控性保证 $\mathcal C(A,B)$ 的列张成 $\mathbb R^n$。任意 $w\in\mathbb R^n$
都能写成 $w=\mathcal C(A,B)\eta$，于是
$Xw=X\mathcal C(A,B)\eta=0$。即 $X$ 把每个向量都映为零，只能是零矩阵。
\end{proofstep}

\pagebreak
\begin{proofstep}{充分性：由不能控性构造一个非零反例矩阵}
假定题述矩阵条件成立。用反证法，设 $(A,B)$ 不能控。
由 PBH 判据，存在 $\lambda\in\mathbb C$ 和非零复行向量 $q$，使
\[
qA=\lambda q,\qquad qB=0.
\]
要得到与 $A$ 交换的矩阵，可以把同一特征值对应的左、右特征向量作外积。
因为 $\lambda$ 是 $A$ 的特征值，$A-\lambda I$ 奇异，故存在非零列向量 $v$ 满足
$Av=\lambda v$。令 $X_{\mathbb C}=vq$。

该矩阵一定非零：取 $v_i\ne0$、$q_j\ne0$，则其 $(i,j)$ 元素为
$v_iq_j\ne0$。逐项检查题述两条件：
\[
AX_{\mathbb C}=(Av)q=\lambda vq,
\qquad X_{\mathbb C}A=v(qA)=\lambda vq,
\]
\[
X_{\mathbb C}B=v(qB)=0.
\]
注意这里不需要 $qv\ne0$；需要非零的是外积 $vq$，而不是内积 $qv$。
\end{proofstep}

\begin{proofstep}{把复矩阵反例转换为实矩阵反例}
若题述 $X$ 限定为实矩阵，写 $X_{\mathbb C}=U+\mathrm i V$，其中 $U,V$ 为实矩阵。
因 $A,B$ 均为实矩阵，把已证等式分别取实部、虚部，得到
\[
AU=UA,\quad UB=0,\qquad AV=VA,\quad VB=0.
\]
而 $X_{\mathbb C}\ne0$ 意味着 $U,V$ 至少一个非零，取该非零者作 $X$，
即与题设的“只能有 $X=0$”矛盾。因此 $(A,B)$ 必须能控，充分性得证。
\end{proofstep}
\end{solution}

\begin{problem}[原题 14]
设 $A\in\mathbb R^{n\times n}$、$B\in\mathbb R^{n\times m}$、$C\in\mathbb R^{p\times n}$。证明：增广系统
\[
A_e=\begin{bmatrix}A&0\\C&0\end{bmatrix},\qquad
B_e=\begin{bmatrix}B\\0\end{bmatrix}
\]
能控的充要条件为 $(A,B)$ 能控，且 $\begin{bmatrix}A&B\\C&0\end{bmatrix}$ 行满秩。
\end{problem}
\begin{solution}
\begin{proofstep}{写出增广系统的 PBH 矩阵}
原状态维数为 $n$，增广后为 $n+p$，因此要证明的秩条件是
$\operatorname{rank}[\lambda I_{n+p}-A_e\ B_e]=n+p$ 对所有复数 $\lambda$ 成立。
把它按原状态、新增状态和输入三组列分块，得到
\[
M(\lambda)=\begin{bmatrix}\lambda I_n-A&0&B\\-C&\lambda I_p&0\end{bmatrix}.
\]
这里中间的 $\lambda I_p$ 在 $\lambda=0$ 时失去可逆性，因此必须把零点与非零点分开检查。
\end{proofstep}

\begin{proofstep}{求出非零点与零点处的秩}
当 $\lambda\ne0$ 时，对第一组列加上“第二组列乘以 $\lambda^{-1}C$”。
这是可逆列变换，结果为
\[
M(\lambda)\ \longrightarrow\
\begin{bmatrix}\lambda I_n-A&0&B\\0&\lambda I_p&0\end{bmatrix}.
\]
下方的 $p$ 行由中间可逆块提供 $p$ 个独立方向，其余秩由上方决定，故
\[
\operatorname{rank}M(\lambda)
=p+\operatorname{rank}[\lambda I_n-A\ \ B].
\]
当 $\lambda=0$ 时，中间 $p$ 列全零，删去零列得 $[-A\ B;-C\ 0]$。
将前 $n$ 列同时乘以 $-1$，得到
\[
\operatorname{rank}M(0)
=\operatorname{rank}\begin{bmatrix}A&B\\C&0\end{bmatrix}.
\]
\end{proofstep}

\begin{proofstep}{充分性：用两个假设覆盖全部复数点}
假定 $(A,B)$ 能控且题述块矩阵行满秩。
对任意 $\lambda\ne0$，原系统能控保证
$\operatorname{rank}[\lambda I_n-A\ B]=n$，故 $\operatorname{rank}M(\lambda)=n+p$。
在 $\lambda=0$ 处，块矩阵行满秩又保证 $\operatorname{rank}M(0)=n+p$。
因此增广系统满足全部 PBH 秩条件，必然能控。
\end{proofstep}

\begin{proofstep}{必要性：从增广系统能控推出两个条件}
假定增广系统能控，则 $M(\lambda)$ 对每个 $\lambda$ 都有满行秩 $n+p$。
取 $\lambda=0$，立即得到题述块矩阵行满秩。
取任意 $\lambda\ne0$，则由第二步的秩公式得到原系统在该点的 PBH 秩为 $n$。

还须补查原系统的 $\lambda=0$ 点。
由于块矩阵的全部 $n+p$ 行线性无关，其中的前 $n$ 行也线性无关，故
\[
\operatorname{rank}[A\ B]=n
\quad\Longleftrightarrow\quad
\operatorname{rank}[-A\ B]=n.
\]
这正是零点处的 PBH 条件。原系统在所有复数点的 PBH 秩均为 $n$，所以 $(A,B)$ 能控。
\end{proofstep}
\end{solution}

\begin{problem}[原题 15]
证明：若 $(A,B)$ 能控，则带已知函数 $f$ 的系统
\[
\dot x=Ax+Bu+f(t)
\]
仍能控。
\end{problem}
\begin{solution}
\begin{proofstep}{把已知外力的作用从终态条件中分离出来}
设已知 $f(t)$ 在有限区间可积，使状态解存在。
为了证明能控，只需对任意初态构造一个把它导到原点的输入。
下面进一步证明：给定任意 $T>0$、初态 $x_0$ 和目标 $x_f$，都能找到输入使 $x(T)=x_f$。

由变常数公式，初始条件、控制输入及外力分别贡献三部分：
\[
x(T)=e^{AT}x_0+\int_0^T e^{A(T-t)}Bu(t)\,dt
+\underbrace{\int_0^T e^{A(T-t)}f(t)\,dt}_{d_f}.
\]
由于 $f$ 已知，$d_f$ 也是已知向量。把它与自由响应一起移到等号另一边，
所求控制必须满足
\[
\int_0^T e^{A(T-t)}Bu(t)\,dt=r,
\qquad r=x_f-e^{AT}x_0-d_f.
\]
因此问题变成：控制输入能否产生任意指定的位移 $r$。
\end{proofstep}

\begin{proofstep}{证明用于构造输入的 Gram 矩阵可逆}
定义终态形式的能控 Gram 矩阵
\[
W_T=\int_0^T e^{A(T-t)}BB^{\mathsf T}e^{A^{\mathsf T}(T-t)}\,dt.
\]
这是实对称矩阵。对任意实列向量 $q$，令 $s=T-t$，有
\[
q^{\mathsf T}W_Tq=\int_0^T
\left\|B^{\mathsf T}e^{A^{\mathsf T}s}q\right\|^2\,ds\ge0.
\]
若等号成立，由被积函数连续且非负，知
$B^{\mathsf T}e^{A^{\mathsf T}s}q=0$ 在整个区间恒成立。
将其在 $s=0$ 处逐次求导，得到
\[
B^{\mathsf T}(A^{\mathsf T})^kq=0\quad(k=0,\ldots,n-1),
\quad\text{即}\quad q^{\mathsf T}[B\ AB\ \cdots\ A^{n-1}B]=0.
\]
由于 $(A,B)$ 能控，该矩阵行满秩，只能有 $q=0$。
故对非零 $q$ 二次型严格为正，即 $W_T>0$，从而 $W_T^{-1}$ 存在。
\end{proofstep}

\begin{proofstep}{构造控制并代回检查终态}
根据 Gram 矩阵的被积形式，选择
\[
u(t)=B^{\mathsf T}e^{A^{\mathsf T}(T-t)}W_T^{-1}r.
\]
这样代入终态积分时，前两项乘积正好成为 $W_T$ 的被积函数：
\begin{align*}
\int_0^T e^{A(T-t)}Bu(t)\,dt
&=\left(\int_0^T e^{A(T-t)}BB^{\mathsf T}
e^{A^{\mathsf T}(T-t)}\,dt\right)W_T^{-1}r\\
&=W_TW_T^{-1}r=r.
\end{align*}
因此可用的显式控制为
\[
\boxed{u(t)=B^{\mathsf T}e^{A^{\mathsf T}(T-t)}W_T^{-1}
\left(x_f-e^{AT}x_0-\int_0^T e^{A(T-\tau)}f(\tau)\,d\tau\right).}
\]
相应终态满足 $x(T)=e^{AT}x_0+r+d_f=x_f$。
特别取 $x_f=0$，就能将任意初态导到原点，故带已知外力的系统仍能控。
已知外力只改变目标位移 $r$，不改变控制输入所能产生位移的线性空间。
\end{proofstep}
\end{solution}

\begin{problem}[原题 20]
设 $(A,B)$ 不完全能控，$T=[T_1\ T_2]$ 非奇异，$T_1$ 的列构成 $X_C$ 的一组基。令
\[
x=Tz,\qquad \widehat A=T^{-1}AT,\quad\widehat B=T^{-1}B.
\]
（1）证明能控子空间随该坐标变换对应；（2）举例说明不能控子空间一般不按同一坐标变换对应；（3）若 $T_2$ 的列构成 $X_{NC}$ 的基，证明不能控子空间也按该坐标变换对应。
\end{problem}
\begin{solution}
\begin{proofstep}{明确坐标映射与子空间定义}
由于本题写的是 $x=Tz$，同一个状态从旧坐标变为新坐标时，实际使用的是
$z=T^{-1}x$。因此要检查的是 $T^{-1}$ 对旧坐标子空间的作用。

记 $\mathcal C(A,B)=[B\ AB\ \cdots\ A^{n-1}B]$。按教材第2.1节（第31页）的定义，
\[
X_C(A,B)=\operatorname{im}\mathcal C(A,B),\qquad
X_{NC}(A,B)=X_C(A,B)^\perp.
\]
这里 $X_{NC}$ 是指定的正交补空间，不是所有不能控状态组成的集合。
这一区别是第（2）问的关键。
\end{proofstep}

\begin{proofstep}{第（1）问：逐列计算能控性矩阵的变换}
对 $k=1$，相邻的 $T,T^{-1}$ 抵消，得到
$\widehat A\widehat B=(T^{-1}AT)(T^{-1}B)=T^{-1}AB$。
同样对任意 $k\ge0$，有
\[
\widehat A^k\widehat B=T^{-1}A^kB,
\quad \mathcal C(\widehat A,\widehat B)=T^{-1}\mathcal C(A,B).
\]
因此新矩阵每一列都等于旧矩阵相应列左乘 $T^{-1}$，其列空间满足
\[
\boxed{X_C(\widehat A,\widehat B)=T^{-1}X_C(A,B),}
\]
即 $x\in X_C(A,B)$ 当且仅当 $z=T^{-1}x$ 属于新系统的能控子空间。
进一步，设 $r=\dim X_C$，由 $T=[T_1\ T_2]$ 可得
\[
T^{-1}T_1=\begin{bmatrix}I_r\\0\end{bmatrix},\qquad
X_C(\widehat A,\widehat B)=\operatorname{span}\{e_1,\ldots,e_r\}.
\]
也就是说，选择能控子空间的一组基作为 $T$ 的前 $r$ 列后，能控部分恰占据前 $r$ 个新坐标。
\end{proofstep}

\begin{proofstep}{第（2）问：构造并核验一个二维反例}
为了让反例计算简单，取输入只作用于第一坐标的系统
\[
A=\begin{bmatrix}0&0\\0&1\end{bmatrix},\quad
B=\begin{bmatrix}1\\0\end{bmatrix},\quad
T=\begin{bmatrix}1&1\\0&1\end{bmatrix},\quad
T^{-1}=\begin{bmatrix}1&-1\\0&1\end{bmatrix}.
\]
$T$ 的第一列为 $e_1$，满足题设要求，第二列 $[1\ 1]^{\mathsf T}$ 则只要求与它线性无关。
由于 $AB=0$，原能控性矩阵为
\[
\mathcal C(A,B)=\begin{bmatrix}1&0\\0&0\end{bmatrix},
\quad X_C=\operatorname{span}\{e_1\},\quad X_{NC}=\operatorname{span}\{e_2\}.
\]
计算变换后的系数，得到
\[
\widehat A=\begin{bmatrix}0&-1\\0&1\end{bmatrix},\qquad
\widehat B=\begin{bmatrix}1\\0\end{bmatrix}.
\]
仍有 $\widehat A\widehat B=0$，所以新系统的能控性矩阵与原矩阵相同，进而
$\widehat X_C=\operatorname{span}\{e_1\}$、$\widehat X_{NC}=\operatorname{span}\{e_2\}$。
但把原不能控子空间中的基向量按坐标映射变换，得到
\[
T^{-1}X_{NC}=\operatorname{span}\left\{\begin{bmatrix}-1\\1\end{bmatrix}\right\}
\ne\widehat X_{NC}.
\]
因为 $[-1\ 1]^{\mathsf T}$ 与新能控方向 $e_1$ 的内积为 $-1$，它不属于新能控空间的正交补。
这就完成了所需反例。
\end{proofstep}

\begin{proofstep}{解释反例：正交补遵循另一种变换规律}
一般非奇异坐标变换不保持欧氏内积，因此不能直接把“取正交补”与“左乘 $T^{-1}$”交换。
为了求正确关系，任取新坐标向量 $w$，则
\begin{align*}
w\in\widehat X_{NC}
&\iff w^{\mathsf T}T^{-1}v=0\quad\text{对所有 }v\in X_C\\
&\iff (T^{-\mathsf T}w)^{\mathsf T}v=0\quad\text{对所有 }v\in X_C\\
&\iff T^{-\mathsf T}w\in X_{NC}.
\end{align*}
故一般成立的是
\[
\boxed{X_{NC}(\widehat A,\widehat B)=T^{\mathsf T}X_{NC}(A,B).}
\]
它与第（1）问中的 $T^{-1}X_C$ 通常不同，正是反例产生的原因。
\end{proofstep}

\begin{proofstep}{第（3）问：用额外的基底条件恢复坐标对应}
现在要求 $T_2$ 的列恰好构成 $X_{NC}$ 的一组基，即 $X_{NC}=\operatorname{im}T_2$。
于是
\[
T^{-1}T_2=\begin{bmatrix}0_{r\times(n-r)}\\I_{n-r}\end{bmatrix}.
\]
所以任意 $x=T_2\eta\in X_{NC}$，在新坐标中都写成 $z=[0\ \eta^{\mathsf T}]^{\mathsf T}$。
第（1）问已证 $\widehat X_C$ 是前 $r$ 个坐标轴所张空间，
故其正交补恰为后 $n-r$ 个坐标轴所张空间。因此
\[
\boxed{T^{-1}X_{NC}(A,B)
=\operatorname{im}\begin{bmatrix}0\\I_{n-r}\end{bmatrix}
=X_{NC}(\widehat A,\widehat B).}
\]
这里要求的是 $T_1,T_2$ 张成相互正交的子空间，不要求各组基内部单位正交。
\end{proofstep}
\end{solution}

\begin{problem}[原题 22]
求下列单输入系统的能控规范型及变换矩阵：
\[
\dot x=\begin{bmatrix}-1&-2&-2\\0&-1&1\\1&0&1\end{bmatrix}x
+\begin{bmatrix}2\\0\\1\end{bmatrix}u,
\qquad y=\begin{bmatrix}0&1&1\end{bmatrix}x.
\]
\end{problem}
\begin{solution}
\begin{proofstep}{先检查系统能控性}
将题面矩阵分别记作 $A,b,c$。为了构造三阶能控规范型，先确认
$b,Ab,A^2b$ 三个方向线性无关。逐次左乘 $A$ 得
\[
Ab=\begin{bmatrix}-4\\1\\3\end{bmatrix},\qquad
A^2b=\begin{bmatrix}-4\\2\\-1\end{bmatrix},\qquad
\mathcal C=[b\ Ab\ A^2b]=\begin{bmatrix}2&-4&-4\\0&1&2\\1&3&-1\end{bmatrix}.
\]
将第一行减去第三行的2倍，再按第一列展开：
\[
\det\mathcal C=
\det\begin{bmatrix}0&-10&-2\\0&1&2\\1&3&-1\end{bmatrix}
=\det\begin{bmatrix}-10&-2\\1&2\end{bmatrix}=-18\ne0.
\]
所以系统完全能控，确实可以用三个新状态构造能控规范型。
\end{proofstep}

\pagebreak
\begin{proofstep}{计算特征多项式并确定目标形式}
按 $sI-A$ 的第一行展开，有
\begin{align*}
\det(sI-A)
&=\det\begin{bmatrix}s+1&2&2\\0&s+1&-1\\-1&0&s-1\end{bmatrix}\\
&=(s+1)^2(s-1)+2+2(s+1)\\
&=s^3+s^2+s+3.
\end{align*}
采用教材第2.4节（第49--51页）的末行伴随型，目标为
\[
A_c=\begin{bmatrix}0&1&0\\0&0&1\\-3&-1&-1\end{bmatrix},\qquad b_c=e_3.
\]
末行依次放特征多项式常数项、一次项、二次项系数的相反数。
\end{proofstep}

\begin{proofstep}{从目标形式反推变换矩阵的三列}
约定 $x=Tz$，因此需要 $AT=TA_c$ 和 $b=Tb_c$。
写 $T=[t_1\ t_2\ t_3]$。因 $b_c=e_3$，首先必须取 $t_3=b$。
再逐列展开 $AT=TA_c$，要求
\[
At_1=-3t_3,\qquad At_2=t_1-t_3,\qquad At_3=t_2-t_3.
\]
从最后一个等式向前求，得到
\[
t_2=At_3+t_3=Ab+b,\qquad
t_1=At_2+t_3=A^2b+Ab+b.
\]
剩下第一个等式由 Cayley--Hamilton 定理自动保证：
\[
At_1=(A^3+A^2+A)b=-3b=-3t_3.
\]
这说明下述列向量的选择来自规范型要求，并非试凑：
\[
T=[A^2b+Ab+b\ \ Ab+b\ \ b]
=\begin{bmatrix}-6&-2&2\\3&1&0\\3&4&1\end{bmatrix},
\qquad \det T=18,
\]
因此 $T$ 非奇异。解 $TT^{-1}=I$ 得
\[
T^{-1}=\begin{bmatrix}
\tfrac1{18}&\tfrac59&-\tfrac19\\
-\tfrac16&-\tfrac23&\tfrac13\\
\tfrac12&1&0
\end{bmatrix}.
\]
\end{proofstep}

\pagebreak
\begin{proofstep}{同时变换状态方程、输入列与输出行}
由 $x=Tz$，原状态方程等价于
$\dot z=T^{-1}ATz+T^{-1}bu$，输出方程等价于 $y=cTz$。
前一步已保证 $T^{-1}AT=A_c$、$T^{-1}b=e_3$，另外
\[
cT=[0\ 1\ 1]T=[3+3\ \ 1+4\ \ 0+1]=[6\ 5\ 1].
\]
因此最终结果为
\[
\boxed{
\dot z=\underbrace{\begin{bmatrix}0&1&0\\0&0&1\\-3&-1&-1\end{bmatrix}}_{T^{-1}AT}z
+\underbrace{\begin{bmatrix}0\\0\\1\end{bmatrix}}_{T^{-1}b}u,
\qquad y=\underbrace{\begin{bmatrix}6&5&1\end{bmatrix}}_{cT}z.}
\]
\end{proofstep}

\begin{proofstep}{用输入输出关系核对规范型}
在零初始条件下，前两行给出 $Z_2=sZ_1$、$Z_3=s^2Z_1$；
第三行给出 $(s^3+s^2+s+3)Z_1=U$。
因此 $Y=(6+5s+s^2)Z_1$，从而
\[
G(s)=\frac{s^2+5s+6}{s^3+s^2+s+3},
\]
与直接计算原系统的 $c(sI-A)^{-1}b$ 相同。
这同时核对了规范型的分母系数和输出行的排列顺序。
\end{proofstep}
\end{solution}
\end{document}
